
 \section{Dynamical weights}
 

\sss{Local weights}
We will now consider a renormalizable $\psi$-ql map
$$  F\equiv (f,i): U' \ra U $$
with little  Julia sets $ J_n =\di K_n $, $n=0, 1, \dots, p-1$, and their
preimages $J_n' = \di K_n' = -J_n$, $n=1,\dots, p-1$.
(As usual, we identify $J_n/ K_n$  in $U'$ with $J_n/K_n$ in $U$.)
% We let $\JJ= \cup J_n$ and $\JJ' = \cup J_n'$. 

The following lemma shows that 
% Let us consider
all the local weights $\WW_n^\KK = \WW (J_n) $ are $2-$comparable:  

   
\begin{lem}\label{comparison of dyn loc w}
We have:
$$
       \WW_1^\KK \leq \dots \leq \WW_{p-1}^\KK  \leq \WW_0^\KK \leq 2\WW_1^\KK
$$
\end{lem}

\bignote{Prepare the ransformation Rule!} 

We let $\WW_\loc^\KK =  \WW_0^\KK \equiv \max \WW_k^\KK $.

The local weight $\WW(\di U)= (1/2) \, \WW(\di U') $
is called $\WW^\ver$.

\sss{Local fluxes}
%Given an $1\leq n\leq p-1$, 
%let us consider  the domain $ D_n $
%%% the group of all little Julia set 
%consisiting of  two adjacent sectors, $n$th and $(n+1)$st.
%%% and all local edges connecting these Julia sets.  
%%%Let $W_\loc^\perp (n)$ be the total flux passing through these edges,
%%%plus the sum of all local weights through the little Julia sets
%%%involved (of course, in the locally bounded case,
%%%we can take only one of them). 
%Let us define the corresponding {\em local flux} as   $W_\loc^\perp (n) := \WW(D_n) $.
%


For $0\leq n\leq p-1$, 
let us define {\em local flux}  $\WW_n^\Hub $
as the sum of all $\WW(K_m)$ and $\WW^\perp(G_i)$ comprising the
path $\gamma_n$ in the Hubbard tree $\Hub$.

\bignote{Need to extend the above discussion of fluxes to surfaces with
  punctures (at $\alpha$ and $\alpha'$).} 
  
The paths $\gamma_n$ form a Markov tiling of $\Hub$,
so we can consider  the corresponding $p\times p$  {\em Hubbard matrix} $A_\Hub$. 
Applying the Transformation Rules  (Lemma \ref{generaltransform rules})
and the Parallel Law (Lemma \ref{transform rule for concatenation}), we obtain

\begin{lem}
      $\displaystyle {\WW_n^\Hub \leq \sum_{A_{nm} =1} \WW_m^\Hub } $
\end{lem}  

Putting this together with Lemma \ref{bounded deg}, we conclude:

\begin{lem}\label{loc fluxes} 
  All local fluxes $W_n^\Hub  $ are comparable,
  with a constant  depending only on $\bb$.
%  the combinatorial bounds,  i.e., on  the number of eyes). 
Hence
$
     \WW(\SAS) \asymp p \WW_\loc^\Hub.  
$
\end{lem}

We let $\WW^\Hub_\loc : = \max \WW^\Hub_n$.
Note that by definition $\WW^\Hub_\loc \geq \WW^\KK_\loc$
but {\em a priori} they can be incomparable.    
We will see below that this is not the case in our situation. 

\msk

\bignote{Now $\WW^\Ray_n$ is put on the same footing.
Call them ``external'' fluxes? (while  the fluxes through $\Hub$
``internal''?
Use notation $\WW_n$, $\WW_n^{\mathrm in}$, and $\WW_n^\ext$?} 

\sss{Segment diagram weight} 
% \sss{Bound on the segment weights}
% \sss{Self-intersecting segments}

We call a segment $\si$  in the translation region
``self-intersecting''  if 
its translation $\si'$  by $-1$ crosses itself
(amounting to self-intersection on the quotient  cylinder.) 

\begin{lem}\label{self-intersecting segments} 
  The total weight of self-interseting segments in  the translation
  region is $O(p+\WW_\loc^\Hub)$.   \note{Check!!} 
\end{lem}

\begin{proof}
For a self-intersecting sgment $\si$, the translation $\si'$ cannot
represent the lift $f^*(\si)$, since otherwise tho segments in $U'\sm
\KK'$ would intersect. The alternative is that $f^*(\si)$ is broken.
By Lemma \ref{key estimate}, the total weight of these segments is
$O(p+\WW^\Hub_\loc)$.
\end{proof}

Let us call a segment $\si$ {\em peripheral} if its translation $\si'$
exits the translation region.   \note{define!} 
Since the total weight of these segments is  $O(\WW^\Hub_\loc)$, \note{OK?}
% $O(\WW^\perp_\loc+\WW^\perp(\RR_0)+\WW^\ver)$,
we conclude: 

\begin{cor}
  The total weight of self-intersecting and peripheral segments in the
  translation region is
  $O(p+\WW_\loc^\Hub)$. 
% $O(p+\WW_\loc^\Hub+\WW^\perp(\RR_0)+\WW^\ver)$. 
\end{cor}  

\bignote{The error  term used to contain $\WW^\perp(\RR_0)+\WW^\ver$.
Why?}

The same analysis applies to the complement of the translation region
\note{be more precise!} leading to the following conclusion:

\begin{cor}\label{long and peripheral} 
  The total weight of self-intersecting and peripheral segments   is
   $O(p+\WW_\loc^\Hub)$. 
%  $O(p+\WW_\loc^\Hub+\WW^\perp(\RR_0)+\WW^\ver)$. 
\end{cor}

With this in hands, only short segments in the translation region are left unattended.

\begin{lem}
  The total weight of short segments is $O(p+\WW^\ver)$.
\end{lem}

\begin{proof}
Indeed, any short segment extends to either one of the segments listed
in Corollary \ref{long and peripheral},
or a loop surrounding a little Julia set, 
or  a vertical arc-segment,
or else to a ``snake'' composed of inner and outer short segments.
Note now that
loops around little Julia sets have small weight
(since they are dual to the local fluxes),
while snake's weights are bounded  by one (since they
intersect their translations). The conclusion follows. 
%and their total weight can be bounded as in
%Lemma \ref{self-intersecting segments}.
\end{proof}  

Putting these together, we obtain:

\begin{cor}\label{total segments weight}
  The total weight of the segment diagram in $U\sm \KK$ is
  bounded by  $O(p + \WW^\Hub_\loc + \WW^\ver)$.
\end{cor}

For $\la>1$, let us say that we  enjoy $\la$-{\em ampliphication} if
$ \WW^\ver \geq \la \WW^\KK_\loc$. 

\begin{cor}\label{total segments weight 2}
 For any $\la>1$,  either we enjoy $\la$-ampliphication 
  or else
  $$ \WW(\SSS_{U\sm \KK }) = O_\la (p + \WW^\Hub_\loc ). $$
\end{cor}



\sss{Arc-segment and long arc diagrams} 


\begin{lem}\label{arc-segments and long arcs}
  The total weight of arc-segments and long%
\footnote{in both the translation region and its complement} 
arcs   in  $U\sm \KK$ 
  is  bounded by  the total weight of segments in $U\sm \KK$. 
% We have:
%$$
%    \WW(\AS) + \WW(\AAA^\long) \leq \WW
%$$
\end{lem} 

\begin{proof}
For the same reason as for long segments in $U\sm \KK$,
 all arc-segments and long arcs get broken under the immersion 
$i:  U'\sm \KK \ra U\sm \KK $.  By   Lemma \ref{broken arc-segments},
their total weight is bounded by the total weight of segments in 
$U'\sm \KK$. But the latter is bounded by the total weight of segments
in $U\sm \KK$. 
\end{proof}

Combining this with Corollary \ref{total segments weight 2},
we obtain:

\begin{cor}\label{everything} 
  For any $\la>1$,  either we enjoy $\la$-ampliphication  or else
  $$ \WW(\SSS) + \WW(\AS) + \WW(\AAA^\lo)  + \WW(\AAA^\pe) = O_\la (p + \WW^\Hub_\loc ), $$
where all the diagrams are considered in $U\sm \KK$.   
\end{cor}  

\sss{Conclusions}

Since the total weight $\WW(\SAS)$ of segments, arcs, and arc-segments in $U\sm \KK$
is of order $p \WW_\loc^\Hub$ (by Lemma \ref{loc fluxes}), we conclude: 

\begin{cor}\label{small proportion}
   For any $\la>1$,  either we enjoy $\la$-ampliphication,  or else
 \begin{equation}\label{small proportion estimate}
   \WW(\SSS) + \WW(\AS) + \WW(\AAA^\lo) + \WW(\AAA^\pe)    = o (\WW(\SAS) )\
   {\mathrm{as}} \  p\to \infty, \ % {\mathrm{and}}\
   \WW(\SAS)\to \infty,
\end{equation}    
where all the diagrams are considered in $U\sm \KK$.     
\end{cor}

%\begin{lem}
%  The total weight of segments,  arc-segments, and long arcs in $U\sm \Kfilled$
%is  %  $O(p+\WW^\Hub_ \loc + \WW^\ver)$.
% $O(p + \WW^\ver)$. \note{elaborate!} 
%\end{lem}

In summary, we obtain:

\begin{thm}\label{short arcs}
  For any $\la>1$,  either we enjoy $\la$-ampliphication,  or else
  the elephant's share of the total weight $\WW(\SAS)$ is supprted on
  short arcs of the translation region%
\footnote{which is the same as ``short arcs in the complement of the
  translation region''}
  $($provided $p$ and $\WW(\SAS)$ are sufficiently big$)$. 
\end{thm}  

\begin{proof}
  What is unaccounted in the left hand side of
  (\ref{small proportion estimate}) is the weight of short arcs.
\end{proof}  

\begin{cor}\label{loc weights comparison} 
%   Under the above circumstances, \note{informally speaking} 
% if there is no amplification, then
If the second option of the above Theorem is realized then 
  
\ssk\nin {\rm (i)} 
$\WW_\loc^\Hub\asymp \WW_\loc^\KK$. 

\ssk\nin {\rm (ii)}
For any $\kappa\in (0,1)$ and $\eps>0$,
if the period $p$ and the  local weight $\WW_\loc^\KK$ are sufficiently big
then for any $n \in [1,\kappa \qq] $
at least $ (1-\eps) $-proportion of the local weight $\WW_n^\KK$ is supported on
the two short arcs emanated from $K_n$.  
%
%\ssk\nin {\rm (ii)}
% There exist absolute $t$ and $C$ such that 
%for $n \leq \qq - t$ at least 
%$ \WW_n^\KK - C$ part of the local weight $\WW_n^\KK$ is supported on
%the two short arcs emanated from $K_n$.  
\end{cor}

\bignote{The error term in (ii) used to be additive $O(1)$.
It can be selected as $O(1/p)$.  }

\begin{proof}
(i)   Indeed, $\WW(\AAA^\sh) = O(p\, \WW_\loc^\KK)$
  while $\WW(\SAS) \asymp p\, \WW_\loc^\Hub $, and
  by   the last theorem, $\WW(\AAA^\sh) \sim \WW(\SAS) $. 

  \msk (ii) It follows from Corollary \ref{small proportion}
  that the assertion is true for some $n_0\in [\kappa \qq, \qq)$.
  Pulling the correspondng  short arcs back, we obtain the desired for all $n\leq n_0$. 
\end{proof}  



% We arrive at the final conclusion of this section:
%
%Thus, up to this error, {\em  the whole weight of the canonical arc-segment
%diagram is alighned with the Hubbard tree. } 

